\section{Log and Event Pipeline}

Loki stores logs differently from a traditional log system. Instead of indexing the full content
of each line, it indexes labels and stores the log chunks compressed in object storage
\cite{grafana-loki}. This keeps costs manageable at scale, but it requires some discipline: too
few labels and queries become slow and broad; too many dynamic labels and the index grows in ways
that hurt performance. Getting that balance right is part of the design work, not something that
happens automatically.

The deployment is not a minimal installation. It runs in distributed mode with gateways, separate
ingesters and queriers, object storage, retention policies, rate limits, and caches
\cite{lsst-it-k8s-cookbook}. This is not complexity for its own sake; during a real observing
night, logs are only useful if they can be queried quickly and if the labels are good enough to
narrow the search fast.

The deeper point is that logs and metrics need to be designed together, not as separate systems
that happen to run on the same cluster. A metric can tell you that latency spiked at a specific
time. A log query can show what the service was reporting at that exact moment. A Kubernetes event
can show that a pod was restarting. A network log can show a device dropping packets. None of those
signals alone tells the full story. The operator needs the timeline, and building that timeline is
only possible if the signals share a common labeling scheme and a common sense of time.
